%% kapa: exam=1 round=19 year=2008 period=1 form=A subject=civil_law subject_ko=민법 seq=6 number=6
%% source: Q-Net 감정평가사 기출문제 (cst00309) artlSeq=1001962 file=extracted/1008/1001962/2008년 제19회 감정평가사 1차시험 기출문제/1교시.hwp
%% conversion: kordoc 4.24.51 -> tools/gen1.py (auto); review=proofread-sonnet; flags=none
\begin{question}{6}
민법상 법인에 관한 설명으로 틀린 것은?
\begin{choices}
  \item 재단법인의 정관에서 정관의 변경방법을 정하지 않은 경우에도, 일정한 요건 하에 목적 기타 정관의 규정을 변경할 수 있다.
  \item 법인의 이사회와 감사는 필요기관이 아니다.
  \item 이사의 대표권이 정관에 의하여 제한되어 있다면 이를 등기하지 않아도 제3자에게 대항할 수 있다.
  \item 이사가 여러 명인 경우에 법인의 사무에 관한 대표권은 각자 행사함이 원칙이다.
  \item 이사는 정관이나 총회의 결의로 금지하지 않은 사항에 관하여 타인으로 하여금 특정한 행위를 대리하게 할 수 있다.
\end{choices}
\end{question}
\answer{3}
